```latex
% ============================================================
% Seção: Uso de IA generativa
% Brenda Gabrielle Alves Nascimento, 24.1.4011
% BCC328 - Construção de Compiladores I - Trabalho Prático 1
% Inserir antes da conclusão.
% ============================================================

\section{Uso de IA generativa}
\label{sec:ia}

O desenvolvimento do trabalho contou com o apoio de uma ferramenta
de inteligência artificial generativa, o assistente Claude
(Anthropic), utilizada ao longo das diferentes etapas de implementação.
A ferramenta foi empregada como recurso de apoio à compreensão dos
requisitos, discussão de decisões de projeto, consulta sobre conceitos
e bibliotecas da linguagem Haskell, elaboração de trechos de código
e investigação de problemas identificados durante a compilação e os
testes.

A interação ocorreu de forma incremental. As solicitações eram
formuladas a partir dos requisitos do trabalho, das decisões de
implementação e dos problemas encontrados durante o desenvolvimento.
As respostas obtidas eram analisadas e comparadas com a especificação,
a documentação das bibliotecas e o material disponibilizado na
disciplina. Quando necessário, as sugestões eram corrigidas, adaptadas
ou descartadas antes de serem incorporadas ao projeto.

A ferramenta também auxiliou na elaboração de módulos e trechos de
código. Entretanto, sua utilização não substituiu a análise dos
requisitos nem a avaliação das soluções adotadas. As decisões de
projeto, a escolha entre alternativas, a integração dos componentes
e a validação do comportamento do compilador permaneceram sob
responsabilidade da autora.

Como parte do processo de verificação, foi adotado o ciclo de
desenvolvimento incremental: definir o problema ou a decisão de
projeto, consultar a ferramenta quando pertinente, analisar a solução
proposta, implementar ou adaptar o código, compilar com
\texttt{-Wall -Werror}, executar os testes e registrar as alterações
aprovadas no Git. Dessa forma, os resultados sugeridos pela IA eram
tratados como propostas a serem verificadas, e não como soluções
automaticamente corretas.

\subsection{Planejamento e organização do desenvolvimento}

\textbf{Contexto da interação.} No início do trabalho, foi necessário
organizar as etapas de implementação a partir do enunciado e da
especificação, considerando as dependências entre os componentes
do compilador.

\textbf{Apoio da IA.} Foi solicitado ao assistente um plano de
desenvolvimento incremental, contemplando desde a configuração do
ambiente até a elaboração do relatório e a validação final.

\textbf{Análise e decisões.} A ferramenta propôs um plano dividido
em dez etapas, com uma sequência baseada nas dependências entre
as fases do compilador. O plano foi utilizado como referência para
organizar o trabalho, mantendo-se a possibilidade de ajustar a ordem
e o escopo das atividades conforme os problemas encontrados.

\textbf{Validação.} Durante a configuração do ambiente, foi
identificada uma incompatibilidade entre a versão do GHC 8.8.4
disponibilizada inicialmente pela distribuição e a restrição
\texttt{base \^{}>=4.18.2.1} exigida pelo esqueleto da disciplina.
A instalação foi refeita utilizando GHCup, com o GHC 9.6.7,
versão compatível com o ambiente oficial da disciplina.

\subsection{Consulta e alinhamento com o material da disciplina}

\textbf{Contexto da interação.} Para manter a organização do projeto
compatível com os exemplos utilizados na disciplina, foi realizada
uma consulta ao repositório \texttt{bcc328-2026-2}, disponibilizado
pelo professor.

\textbf{Apoio da IA.} O assistente auxiliou na análise da estrutura
do repositório e na identificação de convenções presentes nos
exemplos, especialmente a organização dos diretórios
\texttt{Frontend/Lexer}, \texttt{Frontend/Parser},
\texttt{Frontend/Syntax} e \texttt{Frontend/Pretty}. Também foi
observada a organização do exemplo \texttt{Exp}, que mantém
implementações distintas de analisador léxico compartilhando
um único arquivo \texttt{Token.hs}.

\textbf{Análise e decisões.} A estrutura do projeto foi reorganizada
para seguir essas convenções. A classe \texttt{Pretty}, baseada em
\texttt{Text.PrettyPrint.HughesPJ}, foi adotada a partir do arquivo
\texttt{Utils/Pretty.hs} do repositório da disciplina.

Durante essa etapa, duas informações fornecidas pelo assistente
precisaram ser revistas. A primeira envolvia o atributo
\texttt{pkgs.haskell.compiler.ghc966}, utilizado no
\texttt{flake.nix} sugerido, que não estava disponível em
\emph{nixpkgs}; a consulta ao arquivo do professor indicou
\texttt{ghc96} como atributo adequado. A segunda dizia respeito
à possibilidade de o Alex manter estado entre linhas. A leitura
do arquivo \texttt{LineLexer.x} mostrou que o wrapper
\texttt{monadUserState} permite manter estado durante a análise.

A decisão de utilizar um \emph{scanner} \emph{ad-hoc} foi mantida,
mas por outro motivo: uma ação do Alex devolve um único
\emph{token}, enquanto a emissão de múltiplos \texttt{DEDENT}
consecutivos exigiria uma fila de \emph{tokens} pendentes no estado
do analisador.

\textbf{Validação.} As convenções adotadas foram comparadas com
os arquivos do repositório da disciplina, e as informações
divergentes foram corrigidas antes de serem utilizadas no projeto.

\subsection{Apoio na implementação do analisador léxico e da camada de indentação}

\textbf{Contexto da interação.} A implementação do analisador léxico
exigia definir as estruturas de posição, lexemas e tokens, além
de tratar as regras específicas de indentação da linguagem.

\textbf{Apoio da IA.} O assistente auxiliou na elaboração dos tipos
\texttt{Pos}, \texttt{Lexeme} e \texttt{Token}, do módulo
\texttt{MiniPythonLexer.hs} e da camada de indentação
\texttt{Indent.hs}. A implementação da indentação foi discutida
em etapas, inicialmente com a pilha de níveis e, depois, com
as exceções descritas na Seção 4.5 da especificação.

\textbf{Análise e decisões.} As palavras \texttt{True} e
\texttt{False} foram retiradas da lista de palavras reservadas
e passaram a ser reconhecidas como \texttt{TBoolLit}, evitando
duas representações para o mesmo valor. Posteriormente, foi
necessário derivar \texttt{Ord} em \texttt{Lexeme} e
\texttt{Token}, requisito da instância \texttt{Stream} do
Megaparsec.

A implementação da camada de indentação em etapas permitiu
identificar um problema na primeira versão: um programa com
linha em branco no meio de um bloco produzia sequências
indevidas de \texttt{DEDENT} e \texttt{INDENT}. Esse resultado
indicou a necessidade de considerar o tratamento previsto no
item 4 da Seção 4.5 antes de finalizar a implementação.

\textbf{Validação.} Os módulos foram compilados e testados com
programas que exercitavam as regras de indentação, incluindo
casos com linhas em branco e encerramento de blocos. Os
comportamentos divergentes foram utilizados para orientar
as correções.

\subsection{Modelagem da AST e representação textual}

\textbf{Contexto da interação.} A construção da árvore sintática
abstrata (AST) deveria seguir a Seção 3 do enunciado, armazenando
a posição de cada nó. Também era necessário implementar uma
representação textual própria, sem depender de instâncias
\texttt{Show} geradas automaticamente.

\textbf{Apoio da IA.} O assistente auxiliou na elaboração dos
módulos \texttt{MiniPythonSyntax.hs} e
\texttt{MiniPythonPretty.hs}, incluindo a organização dos tipos
da AST e a impressão de expressões considerando os níveis
de precedência.

\textbf{Análise e decisões.} A primeira versão apresentou
problemas de compilação relacionados a instâncias órfãs.
Como \texttt{-Wall} inclui o aviso \texttt{-Worphans}, foi
necessário tratar explicitamente essa condição. A solução
adotada utilizou \texttt{\{-\# OPTIONS\_GHC -Wno-orphans \#-\}}
acompanhada de comentário justificando a separação entre
a classe, os tipos e as instâncias de \texttt{Pretty}. Essa
abordagem também aparece no arquivo \texttt{LineLexer.x}
do professor.

Outros ajustes envolveram o acesso a um registro, em que
\texttt{case clsSuper of} foi corrigido para
\texttt{case clsSuper cd of}, e a composição da saída textual:
o uso de \texttt{punctuate (text "")} com \texttt{vcat}
não produzia a linha em branco esperada, sendo substituído
por \texttt{\$+\$}.

\textbf{Validação.} A AST e o impressor foram verificados por
meio da compilação e da comparação da representação gerada
com a estrutura esperada para os programas de teste.

\subsection{Discussão da gramática e implementação do analisador descendente recursivo}

\textbf{Contexto da interação.} Foi necessário definir uma
gramática livre de contexto (GLC) não ambígua e implementar
um analisador descendente recursivo organizado em módulos.

\textbf{Apoio da IA.} O assistente auxiliou na discussão das
produções da gramática e na elaboração de uma primeira versão
de \texttt{docs/gramatica.md}, além de apoiar a organização
do analisador nos módulos \texttt{Core.hs}, \texttt{Expr.hs}
e \texttt{RecursiveDescent.hs}.

\textbf{Análise e decisões.} A primeira versão da gramática
precisou ser revista porque a produção \texttt{Alvo} permitia
chamadas de função como alvo de atribuição, deixando essa
restrição apenas em uma explicação textual. A produção foi
reescrita para expressar a restrição diretamente na gramática.

Na implementação do analisador, os quatro programas de teste
falharam inicialmente em expressões como \texttt{int(input())}.
A investigação mostrou que \texttt{int}, \texttt{float},
\texttt{str} e \texttt{bool} são palavras reservadas de tipo,
mas também funcionam como conversões primitivas na Seção 8
da especificação. O não-terminal \texttt{Atomo} não tratava
esses identificadores em posição de expressão, sendo necessário
ajustar a análise sintática.

Também foi revista a convenção de posição dos nós. A primeira
versão utilizava a posição do operador em \texttt{EBin}; por
isso, um erro em \texttt{1 + 2 = 3} apontava para o símbolo
\texttt{+}. Foi adotada a convenção de que a posição de um nó
corresponde ao início da construção no texto-fonte, o que
também favorece a localização de erros em etapas posteriores,
como a verificação de tipos.

\textbf{Validação.} A gramática e o analisador foram avaliados
com os programas de teste da especificação, incluindo casos
válidos e inválidos. As falhas observadas orientaram as
alterações nas produções e nas regras de análise.

\subsection{Comparação de alternativas e implementação do analisador PEG}

\textbf{Contexto da interação.} A implementação do analisador PEG
exigia a utilização da biblioteca \texttt{typed-peg}. Diferentemente
do analisador descendente recursivo, que utiliza uma lista de tokens
produzida pelo analisador léxico, o analisador PEG realiza a análise
diretamente sobre o texto-fonte.

\textbf{Apoio da IA.} O assistente auxiliou na compreensão da
estrutura da biblioteca, na elaboração das regras da gramática PEG
e na implementação do módulo \texttt{PEG.hs}, utilizando o
quasi-quoter \texttt{pegGrammar}. Também contribuiu na investigação
de problemas relacionados ao reconhecimento das construções da
linguagem e à localização dos erros no código-fonte.

\textbf{Análise e decisões.} Durante o desenvolvimento, foi
avaliado o uso da implementação PEG disponibilizada na disciplina
(\texttt{workspace/src/PEG/}). Após a análise de sua estrutura
e dos recursos utilizados, optou-se pela biblioteca
\texttt{typed-peg}, alternativa permitida pelo enunciado.

A implementação exigiu ajustes na gramática para contemplar
o reconhecimento das construções da linguagem diretamente no
texto-fonte, incluindo o tratamento da indentação. Também foi
necessário considerar as informações de posição disponibilizadas
pela biblioteca para produzir mensagens de erro adequadas.

A utilização do \texttt{typed-peg} foi mantida na versão final
do projeto. É importante destacar que o Megaparsec foi utilizado
no analisador descendente recursivo, e não no PEG. Essa separação
permitiu manter as duas implementações sintáticas independentes,
conforme a organização proposta para o trabalho.

\textbf{Validação.} O analisador PEG foi compilado e executado
com os programas de teste da especificação. Seus resultados
foram comparados com os produzidos pelos outros analisadores,
verificando o reconhecimento das construções da linguagem e
a consistência dos resultados.

\subsection{Implementação e depuração do analisador LALR}

\textbf{Contexto da interação.} A implementação do analisador
LALR com Happy deveria declarar a precedência por meio de
diretivas, sem estratificar as produções de expressão,
conforme a Seção 5 do enunciado.

\textbf{Apoio da IA.} O assistente auxiliou na elaboração
inicial do arquivo \texttt{Grammar.y} e na investigação
dos conflitos reportados pelo gerador Happy.

\textbf{Análise e decisões.} A primeira versão apresentou
16 conflitos \emph{shift/reduce}. A investigação utilizou
o arquivo de informações gerado pelo comando
\texttt{happy --info}, que indicou que os conflitos estavam
concentrados em um único estado e relacionados à produção
do corpo da \emph{lambda}.

A primeira tentativa de resolução sugerida, por meio da
declaração \texttt{\%right 'lambda'}, não produziu efeito.
A análise mostrou que, no Happy, a precedência de uma
produção é herdada de seu último terminal, que naquela
regra era o símbolo \texttt{:}. A solução adotada foi
introduzir o pseudo-\emph{token} \texttt{LAMBDA}, com a
menor precedência da tabela, e utilizar
\texttt{\%prec LAMBDA} para sobrescrever a precedência
da produção. Com essa alteração, os conflitos foram
eliminados.

Também foram identificadas dependências ausentes na
configuração do projeto: \texttt{array}, exigida pelo
código gerado, e \texttt{build-tool-depends: happy:happy}
no arquivo \texttt{.cabal}.

\textbf{Validação.} A saída de diagnóstico do Happy foi
utilizada para localizar os conflitos, e a geração do
analisador foi repetida após as alterações, até que
não fossem mais reportados conflitos.

\subsection{Apoio na elaboração e validação dos testes}

\textbf{Contexto da interação.} Era necessário construir
programas de teste com base nos exemplos da Seção 9 da
especificação e verificar se os três analisadores
apresentavam resultados equivalentes.

\textbf{Apoio da IA.} O assistente auxiliou na elaboração
de nove programas válidos, sete programas inválidos e
dos \emph{scripts} \texttt{testes/rodar.sh} e
\texttt{testes/gerar-saidas.sh}, utilizados na execução
e comparação dos analisadores.

\textbf{Análise e decisões.} A primeira execução do
comparativo apresentou falsos positivos: como a compilação
estava quebrada, os dois lados da comparação retornavam
saídas vazias, que eram interpretadas pelo \texttt{diff}
como iguais. O \emph{script} foi então modificado para
armazenar as saídas em arquivos e verificar se elas
não estavam vazias antes de realizar a comparação.

\textbf{Validação.} Os testes foram executados novamente
após a correção do \emph{script}. A verificação passou
a considerar não apenas a igualdade entre as saídas,
mas também se os programas haviam sido executados
corretamente e produzido resultados.

\subsection{Limitações e avaliação das respostas}

Durante o desenvolvimento, algumas respostas fornecidas
pela ferramenta apresentaram inconsistências ou não
consideraram completamente as particularidades do ambiente
e da especificação da disciplina. Entre os casos
identificados estavam a indicação de um atributo
inexistente na configuração do Nix, uma interpretação
incorreta sobre o funcionamento do Alex e uma tentativa
inicial de resolver os conflitos de precedência do Happy
que não produziu o efeito esperado.

Essas situações evidenciaram a importância de avaliar
criticamente as sugestões recebidas. Para isso, foram
consultados o material disponibilizado pelo professor,
a documentação das bibliotecas e os resultados obtidos
com a compilação e os testes. As sugestões que não
atendiam aos requisitos foram corrigidas, adaptadas
ou descartadas.

A experiência mostrou que a IA foi útil como apoio
à investigação e à implementação, especialmente para
explorar alternativas e acelerar a identificação de
possíveis causas de erros. Entretanto, a correção das
soluções dependia de sua verificação no contexto real
do projeto.

\subsection{Síntese}

A IA generativa foi utilizada como assistente técnico
durante diferentes etapas do trabalho. Sua contribuição
envolveu a organização do desenvolvimento, a consulta
sobre conceitos e bibliotecas, a discussão de alternativas
de projeto, a elaboração de trechos de código e o apoio
na investigação de falhas.

A participação da ferramenta foi especialmente útil
em tarefas que exigiram integração entre componentes,
como o tratamento de indentação, a construção dos
analisadores sintáticos e a resolução de conflitos
de precedência no Happy. Ao mesmo tempo, as interações
também permitiram identificar limitações nas respostas
e a necessidade de confrontá-las com a especificação
e com o material da disciplina.

O desenvolvimento foi conduzido de forma incremental,
com compilação utilizando \texttt{-Wall -Werror},
execução de testes e registro das alterações no Git.
As sugestões da ferramenta não foram incorporadas
automaticamente: foram analisadas, adaptadas e
validadas de acordo com os requisitos do trabalho.
Assim, a IA atuou como recurso de apoio à implementação
e à resolução de problemas, enquanto as decisões de
projeto, a análise crítica e a validação final do
compilador permaneceram sob responsabilidade da autora.
```